% explainer-source-hash: 30eba2c07fd93bd3e818b2a9942e6f70ace175b93b997cfd1d710afb07bbfc59
% explainer-source-commit: 6a20de2f6ebc03e4d85c9454ba07bb53ed18f9da
% explainer-generated: 2026-10-15
% Regenerate with /explain-all (see WIRING.md). Do not edit this stamp by hand:
% it is how the toolchain knows whether this document still matches the code.
\documentclass[11pt,a4paper]{article}
% preamble.tex - shared preamble for this project's explainer documents.
%
% This file is generated: it is copied in beside the documents on every build, so a
% clone of this repo can rebuild the PDFs, but local edits to it are overwritten.
%
% Usage from a document:
%   \documentclass[11pt,a4paper]{article}
%   % preamble.tex - shared preamble for this project's explainer documents.
%
% This file is generated: it is copied in beside the documents on every build, so a
% clone of this repo can rebuild the PDFs, but local edits to it are overwritten.
%
% Usage from a document:
%   \documentclass[11pt,a4paper]{article}
%   % preamble.tex - shared preamble for this project's explainer documents.
%
% This file is generated: it is copied in beside the documents on every build, so a
% clone of this repo can rebuild the PDFs, but local edits to it are overwritten.
%
% Usage from a document:
%   \documentclass[11pt,a4paper]{article}
%   \input{preamble}
%   \begin{document} ... \end{document}
%
% Build:  pdflatex -interaction=nonstopmode deep-dive.tex   (run it three times, so the
%         table of contents and the TikZ positioning settle)
%
% PACKAGE NOTES. Every package used here was checked present on the machine these
% documents were built on (Debian TeX Live 2023). If you rebuild elsewhere and a build
% fails, these are the likely gaps. Do not add a package without checking first:
%   kpsewhich <pkg>.sty
%
% Deliberately NOT used, because they were absent and the build dies without them:
%   algorithm2e, algorithm, algorithmic, algpseudocode  (texlive-science not installed)
%     -> the \begin{pseudocode} environment defined below stands in for them.
%   lmodern, charter, mathptmx, times, mathpazo, courier, newtx*, tgtermes
%     -> their .sty files RESOLVE but the font files themselves are absent, so the
%        build dies with "I can't find file `bchr8t'" or similar. Verified by
%        test-compiling each one. Only Computer Modern (the default) and helvet
%        actually render. kpsewhich reporting a font it cannot render is the trap:
%        do not "upgrade" the font stack without test-compiling it first.
%   amsmath, amssymb -> not loaded, so \text{} and \blacksquare are unavailable.
%
% A fuller font stack and algorithm2e would come from:
%   sudo apt install texlive-fonts-recommended texlive-science cm-super

\usepackage[T1]{fontenc}
\usepackage[utf8]{inputenc}
% Font: Computer Modern (default). See the note above before changing this.
\usepackage[a4paper,margin=2.4cm,headheight=14pt]{geometry}
% microtype WITHOUT font expansion. This is load-bearing, do not "restore" it:
% cm-super is not installed, so T1 Computer Modern resolves to BITMAP fonts, and
% pdfTeX's auto-expansion needs scalable ones. With expansion on, any document
% long enough to reach a second page dies with
%   "pdfTeX error (font expansion): auto expansion is only possible with scalable fonts"
% Protrusion (the bigger visual win) still works and stays on.
% `sudo apt install cm-super` would lift this; until then, leave expansion off.
\usepackage[protrusion=true,expansion=false]{microtype}
% Kill hyphen ligatures in TYPEWRITER text. This is load-bearing, not cosmetic. Under T1,
% \texttt{--verbose} silently renders as "-verbose", and \texttt{---} as "--", so a
% document can print a command-line flag as the WRONG flag and look perfectly fine doing
% it. Verified by rendering the page and reading the glyphs. lstlisting and \verb were
% always safe; this makes \code, \file and \texttt safe too, which is what prose uses.
% Ordinary prose is unaffected and still ligates, which is what you want there.
\DisableLigatures[-]{family=tt*}
\usepackage{graphicx}
\usepackage{xcolor}
\usepackage{booktabs}
\usepackage{enumitem}
\usepackage{fancyhdr}
\usepackage{titlesec}
\usepackage{tcolorbox}
\usepackage{listings}
\usepackage{float}
\usepackage{caption}
\usepackage{tikz}
\usepackage{forest}
\usepackage{pgfplots}
\pgfplotsset{compat=1.18}
\usetikzlibrary{arrows.meta,positioning,shapes.geometric,fit,backgrounds,calc,decorations.pathreplacing}
\tcbuselibrary{skins,breakable}
\usepackage[hidelinks,bookmarks=true,bookmarksnumbered=true]{hyperref}

% ===== palette ===========================================================
\definecolor{inkblue}{HTML}{1F3A5F}
\definecolor{accent}{HTML}{2E7DAF}
\definecolor{softgrey}{HTML}{F4F6F8}
\definecolor{linegrey}{HTML}{C8D0D8}
\definecolor{codebg}{HTML}{F7F7F4}
\definecolor{warnred}{HTML}{A33B3B}
\definecolor{okgreen}{HTML}{2E7D53}

% ===== headings ==========================================================
\titleformat{\section}{\Large\bfseries\color{inkblue}}{\thesection}{0.7em}{}
\titleformat{\subsection}{\large\bfseries\color{inkblue}}{\thesubsection}{0.6em}{}
\titleformat{\subsubsection}{\normalsize\bfseries\color{accent}}{\thesubsubsection}{0.5em}{}

\pagestyle{fancy}
\fancyhf{}
\fancyhead[L]{\small\color{inkblue}\nouppercase{\leftmark}}
\fancyhead[R]{\small\color{inkblue}\thepage}
\renewcommand{\headrulewidth}{0.4pt}
\renewcommand{\headrule}{\hbox to\headwidth{\color{linegrey}\leaders\hrule height \headrulewidth\hfill}}

% ===== code ==============================================================
\lstdefinestyle{house}{
  backgroundcolor=\color{codebg},
  basicstyle=\ttfamily\footnotesize,
  keywordstyle=\color{inkblue}\bfseries,
  commentstyle=\color{gray}\itshape,
  stringstyle=\color{okgreen},
  numbers=left, numberstyle=\tiny\color{gray}, numbersep=8pt,
  frame=single, rulecolor=\color{linegrey},
  breaklines=true, breakatwhitespace=false,
  showstringspaces=false, tabsize=2,
  captionpos=b, xleftmargin=14pt, framexleftmargin=14pt,
  columns=fullflexible, keepspaces=true,
  upquote=true,  % render ' and " as straight quotes, not curly
  % NON-ASCII INSIDE A LISTING IS A HARD BUILD FAILURE, not a cosmetic issue: listings
  % reads raw bytes, so inputenc dies on the first multi-byte character ("Invalid UTF-8
  % byte sequence") and the whole document fails to build. These mappings cover what
  % realistically turns up in source code: Latin-1 accents, Spanish punctuation, smart
  % quotes picked up from a word processor, the degree and section signs.
  % NOT COVERED, and it cannot be here: Urdu, Arabic and Balochi script. pdflatex has no
  % font for them. Describe or transliterate such text; never paste it into a listing.
  literate=%
    {á}{{\'a}}1 {é}{{\'e}}1 {í}{{\'i}}1 {ó}{{\'o}}1 {ú}{{\'u}}1
    {Á}{{\'A}}1 {É}{{\'E}}1 {Í}{{\'I}}1 {Ó}{{\'O}}1 {Ú}{{\'U}}1
    {à}{{\`a}}1 {è}{{\`e}}1 {ì}{{\`i}}1 {ò}{{\`o}}1 {ù}{{\`u}}1
    {ä}{{\"a}}1 {ë}{{\"e}}1 {ï}{{\"i}}1 {ö}{{\"o}}1 {ü}{{\"u}}1
    {Ä}{{\"A}}1 {Ö}{{\"O}}1 {Ü}{{\"U}}1 {ß}{{\ss}}1
    {ñ}{{\~n}}1 {Ñ}{{\~N}}1 {ç}{{\c c}}1 {Ç}{{\c C}}1
    {¿}{{\textquestiondown}}1 {¡}{{\textexclamdown}}1
    {°}{{\textdegree}}1 {§}{{\S}}1 {£}{{\pounds}}1 {€}{{EUR}}1
    {“}{{\textquotedblleft}}1 {”}{{\textquotedblright}}1
    {‘}{{\textquoteleft}}1 {’}{{\textquoteright}}1
    {–}{{-}}1 {—}{{-}}1 {…}{{...}}1 {→}{{$\rightarrow$}}1 {✓}{{x}}1
}
\lstset{style=house}

% ===== callout boxes =====================================================
% \begin{keyidea}{Title} ... \end{keyidea}   - a concept a beginner must grasp
\newtcolorbox{keyidea}[1]{
  enhanced, breakable, colback=softgrey, colframe=accent,
  boxrule=0.4pt, arc=2pt, left=8pt, right=8pt, top=6pt, bottom=6pt,
  title={\bfseries #1}, coltitle=white, colbacktitle=accent, fonttitle=\small
}

% \begin{jargon}{Term} definition \end{jargon}  - define a term at first use
\newtcolorbox{jargon}[1]{
  enhanced, breakable, colback=white, colframe=linegrey,
  boxrule=0.4pt, arc=2pt, left=8pt, right=8pt, top=6pt, bottom=6pt,
  title={\bfseries Jargon: #1}, coltitle=inkblue, colbacktitle=softgrey, fonttitle=\small
}

% \begin{honest}{Title} ... \end{honest}  - brutally honest finding (rule 18)
\newtcolorbox{honest}[1]{
  enhanced, breakable, colback=white, colframe=warnred,
  boxrule=0.6pt, arc=2pt, left=8pt, right=8pt, top=6pt, bottom=6pt,
  title={\bfseries #1}, coltitle=white, colbacktitle=warnred, fonttitle=\small
}

% \begin{pseudocode}{Caption} ... \end{pseudocode}
% Stands in for algorithm2e, which is NOT installed. Use \Step / \Indent inside.
\newtcolorbox{pseudocode}[1]{
  enhanced, breakable, colback=codebg, colframe=inkblue,
  boxrule=0.4pt, arc=2pt, left=10pt, right=8pt, top=6pt, bottom=6pt,
  fontupper=\ttfamily\small,
  title={\bfseries\rmfamily #1}, coltitle=white, colbacktitle=inkblue, fonttitle=\small
}
\newcommand{\Step}[1]{\par\noindent #1}
\newcommand{\Indent}[1]{\par\noindent\hspace*{2em}#1}

% ===== tikz house styles =================================================
\tikzset{
  box/.style   = {rectangle, rounded corners=2pt, draw=inkblue, fill=softgrey,
                  text=inkblue, align=center, inner sep=6pt, minimum height=9mm,
                  font=\small},
  extbox/.style= {rectangle, rounded corners=2pt, draw=linegrey, fill=white,
                  text=black!70, align=center, inner sep=6pt, minimum height=9mm,
                  font=\small, dashed},
  store/.style = {cylinder, shape border rotate=90, aspect=0.25, draw=accent,
                  fill=white, text=inkblue, align=center, inner sep=5pt,
                  font=\small},
  dec/.style   = {diamond, aspect=1.6, draw=accent, fill=white, text=inkblue,
                  align=center, inner sep=2pt, font=\scriptsize},
  flow/.style  = {-{Stealth[length=2.4mm]}, draw=inkblue, thick},
  dflow/.style = {-{Stealth[length=2.4mm]}, draw=accent, thick, dashed},
  % align=center is load-bearing: without it, a \\ inside an edge label throws
  % the misleading "perhaps a missing \item" error.
  lbl/.style   = {font=\scriptsize\itshape, text=black!65, fill=white,
                  inner sep=2pt, align=center}
}

% forest default for directory trees
\forestset{
  dirtree/.style={
    for tree={
      font=\ttfamily\small, grow'=0, child anchor=west, parent anchor=south,
      anchor=west, calign=first, inner sep=2pt,
      edge path={\noexpand\path [draw, \forestoption{edge}]
        (!u.south west) +(9pt,0) |- (.child anchor) \forestoption{edge label};},
      before typesetting nodes={if n=1{insert before={[,phantom]}}{}},
      fit=band, before computing xy={l=15pt}
    }
  }
}

% ===== misc ==============================================================
\captionsetup{font=small, labelfont={bf,color=inkblue}}
\setlist[itemize]{leftmargin=1.2em, itemsep=2pt, topsep=3pt}
\setlist[enumerate]{leftmargin=1.4em, itemsep=2pt, topsep=3pt}
\newcommand{\file}[1]{\texttt{\small #1}}
\newcommand{\code}[1]{\texttt{\small #1}}
\setcounter{tocdepth}{2}
\setcounter{secnumdepth}{3}

%   \begin{document} ... \end{document}
%
% Build:  pdflatex -interaction=nonstopmode deep-dive.tex   (run it three times, so the
%         table of contents and the TikZ positioning settle)
%
% PACKAGE NOTES. Every package used here was checked present on the machine these
% documents were built on (Debian TeX Live 2023). If you rebuild elsewhere and a build
% fails, these are the likely gaps. Do not add a package without checking first:
%   kpsewhich <pkg>.sty
%
% Deliberately NOT used, because they were absent and the build dies without them:
%   algorithm2e, algorithm, algorithmic, algpseudocode  (texlive-science not installed)
%     -> the \begin{pseudocode} environment defined below stands in for them.
%   lmodern, charter, mathptmx, times, mathpazo, courier, newtx*, tgtermes
%     -> their .sty files RESOLVE but the font files themselves are absent, so the
%        build dies with "I can't find file `bchr8t'" or similar. Verified by
%        test-compiling each one. Only Computer Modern (the default) and helvet
%        actually render. kpsewhich reporting a font it cannot render is the trap:
%        do not "upgrade" the font stack without test-compiling it first.
%   amsmath, amssymb -> not loaded, so \text{} and \blacksquare are unavailable.
%
% A fuller font stack and algorithm2e would come from:
%   sudo apt install texlive-fonts-recommended texlive-science cm-super

\usepackage[T1]{fontenc}
\usepackage[utf8]{inputenc}
% Font: Computer Modern (default). See the note above before changing this.
\usepackage[a4paper,margin=2.4cm,headheight=14pt]{geometry}
% microtype WITHOUT font expansion. This is load-bearing, do not "restore" it:
% cm-super is not installed, so T1 Computer Modern resolves to BITMAP fonts, and
% pdfTeX's auto-expansion needs scalable ones. With expansion on, any document
% long enough to reach a second page dies with
%   "pdfTeX error (font expansion): auto expansion is only possible with scalable fonts"
% Protrusion (the bigger visual win) still works and stays on.
% `sudo apt install cm-super` would lift this; until then, leave expansion off.
\usepackage[protrusion=true,expansion=false]{microtype}
% Kill hyphen ligatures in TYPEWRITER text. This is load-bearing, not cosmetic. Under T1,
% \texttt{--verbose} silently renders as "-verbose", and \texttt{---} as "--", so a
% document can print a command-line flag as the WRONG flag and look perfectly fine doing
% it. Verified by rendering the page and reading the glyphs. lstlisting and \verb were
% always safe; this makes \code, \file and \texttt safe too, which is what prose uses.
% Ordinary prose is unaffected and still ligates, which is what you want there.
\DisableLigatures[-]{family=tt*}
\usepackage{graphicx}
\usepackage{xcolor}
\usepackage{booktabs}
\usepackage{enumitem}
\usepackage{fancyhdr}
\usepackage{titlesec}
\usepackage{tcolorbox}
\usepackage{listings}
\usepackage{float}
\usepackage{caption}
\usepackage{tikz}
\usepackage{forest}
\usepackage{pgfplots}
\pgfplotsset{compat=1.18}
\usetikzlibrary{arrows.meta,positioning,shapes.geometric,fit,backgrounds,calc,decorations.pathreplacing}
\tcbuselibrary{skins,breakable}
\usepackage[hidelinks,bookmarks=true,bookmarksnumbered=true]{hyperref}

% ===== palette ===========================================================
\definecolor{inkblue}{HTML}{1F3A5F}
\definecolor{accent}{HTML}{2E7DAF}
\definecolor{softgrey}{HTML}{F4F6F8}
\definecolor{linegrey}{HTML}{C8D0D8}
\definecolor{codebg}{HTML}{F7F7F4}
\definecolor{warnred}{HTML}{A33B3B}
\definecolor{okgreen}{HTML}{2E7D53}

% ===== headings ==========================================================
\titleformat{\section}{\Large\bfseries\color{inkblue}}{\thesection}{0.7em}{}
\titleformat{\subsection}{\large\bfseries\color{inkblue}}{\thesubsection}{0.6em}{}
\titleformat{\subsubsection}{\normalsize\bfseries\color{accent}}{\thesubsubsection}{0.5em}{}

\pagestyle{fancy}
\fancyhf{}
\fancyhead[L]{\small\color{inkblue}\nouppercase{\leftmark}}
\fancyhead[R]{\small\color{inkblue}\thepage}
\renewcommand{\headrulewidth}{0.4pt}
\renewcommand{\headrule}{\hbox to\headwidth{\color{linegrey}\leaders\hrule height \headrulewidth\hfill}}

% ===== code ==============================================================
\lstdefinestyle{house}{
  backgroundcolor=\color{codebg},
  basicstyle=\ttfamily\footnotesize,
  keywordstyle=\color{inkblue}\bfseries,
  commentstyle=\color{gray}\itshape,
  stringstyle=\color{okgreen},
  numbers=left, numberstyle=\tiny\color{gray}, numbersep=8pt,
  frame=single, rulecolor=\color{linegrey},
  breaklines=true, breakatwhitespace=false,
  showstringspaces=false, tabsize=2,
  captionpos=b, xleftmargin=14pt, framexleftmargin=14pt,
  columns=fullflexible, keepspaces=true,
  upquote=true,  % render ' and " as straight quotes, not curly
  % NON-ASCII INSIDE A LISTING IS A HARD BUILD FAILURE, not a cosmetic issue: listings
  % reads raw bytes, so inputenc dies on the first multi-byte character ("Invalid UTF-8
  % byte sequence") and the whole document fails to build. These mappings cover what
  % realistically turns up in source code: Latin-1 accents, Spanish punctuation, smart
  % quotes picked up from a word processor, the degree and section signs.
  % NOT COVERED, and it cannot be here: Urdu, Arabic and Balochi script. pdflatex has no
  % font for them. Describe or transliterate such text; never paste it into a listing.
  literate=%
    {á}{{\'a}}1 {é}{{\'e}}1 {í}{{\'i}}1 {ó}{{\'o}}1 {ú}{{\'u}}1
    {Á}{{\'A}}1 {É}{{\'E}}1 {Í}{{\'I}}1 {Ó}{{\'O}}1 {Ú}{{\'U}}1
    {à}{{\`a}}1 {è}{{\`e}}1 {ì}{{\`i}}1 {ò}{{\`o}}1 {ù}{{\`u}}1
    {ä}{{\"a}}1 {ë}{{\"e}}1 {ï}{{\"i}}1 {ö}{{\"o}}1 {ü}{{\"u}}1
    {Ä}{{\"A}}1 {Ö}{{\"O}}1 {Ü}{{\"U}}1 {ß}{{\ss}}1
    {ñ}{{\~n}}1 {Ñ}{{\~N}}1 {ç}{{\c c}}1 {Ç}{{\c C}}1
    {¿}{{\textquestiondown}}1 {¡}{{\textexclamdown}}1
    {°}{{\textdegree}}1 {§}{{\S}}1 {£}{{\pounds}}1 {€}{{EUR}}1
    {“}{{\textquotedblleft}}1 {”}{{\textquotedblright}}1
    {‘}{{\textquoteleft}}1 {’}{{\textquoteright}}1
    {–}{{-}}1 {—}{{-}}1 {…}{{...}}1 {→}{{$\rightarrow$}}1 {✓}{{x}}1
}
\lstset{style=house}

% ===== callout boxes =====================================================
% \begin{keyidea}{Title} ... \end{keyidea}   - a concept a beginner must grasp
\newtcolorbox{keyidea}[1]{
  enhanced, breakable, colback=softgrey, colframe=accent,
  boxrule=0.4pt, arc=2pt, left=8pt, right=8pt, top=6pt, bottom=6pt,
  title={\bfseries #1}, coltitle=white, colbacktitle=accent, fonttitle=\small
}

% \begin{jargon}{Term} definition \end{jargon}  - define a term at first use
\newtcolorbox{jargon}[1]{
  enhanced, breakable, colback=white, colframe=linegrey,
  boxrule=0.4pt, arc=2pt, left=8pt, right=8pt, top=6pt, bottom=6pt,
  title={\bfseries Jargon: #1}, coltitle=inkblue, colbacktitle=softgrey, fonttitle=\small
}

% \begin{honest}{Title} ... \end{honest}  - brutally honest finding (rule 18)
\newtcolorbox{honest}[1]{
  enhanced, breakable, colback=white, colframe=warnred,
  boxrule=0.6pt, arc=2pt, left=8pt, right=8pt, top=6pt, bottom=6pt,
  title={\bfseries #1}, coltitle=white, colbacktitle=warnred, fonttitle=\small
}

% \begin{pseudocode}{Caption} ... \end{pseudocode}
% Stands in for algorithm2e, which is NOT installed. Use \Step / \Indent inside.
\newtcolorbox{pseudocode}[1]{
  enhanced, breakable, colback=codebg, colframe=inkblue,
  boxrule=0.4pt, arc=2pt, left=10pt, right=8pt, top=6pt, bottom=6pt,
  fontupper=\ttfamily\small,
  title={\bfseries\rmfamily #1}, coltitle=white, colbacktitle=inkblue, fonttitle=\small
}
\newcommand{\Step}[1]{\par\noindent #1}
\newcommand{\Indent}[1]{\par\noindent\hspace*{2em}#1}

% ===== tikz house styles =================================================
\tikzset{
  box/.style   = {rectangle, rounded corners=2pt, draw=inkblue, fill=softgrey,
                  text=inkblue, align=center, inner sep=6pt, minimum height=9mm,
                  font=\small},
  extbox/.style= {rectangle, rounded corners=2pt, draw=linegrey, fill=white,
                  text=black!70, align=center, inner sep=6pt, minimum height=9mm,
                  font=\small, dashed},
  store/.style = {cylinder, shape border rotate=90, aspect=0.25, draw=accent,
                  fill=white, text=inkblue, align=center, inner sep=5pt,
                  font=\small},
  dec/.style   = {diamond, aspect=1.6, draw=accent, fill=white, text=inkblue,
                  align=center, inner sep=2pt, font=\scriptsize},
  flow/.style  = {-{Stealth[length=2.4mm]}, draw=inkblue, thick},
  dflow/.style = {-{Stealth[length=2.4mm]}, draw=accent, thick, dashed},
  % align=center is load-bearing: without it, a \\ inside an edge label throws
  % the misleading "perhaps a missing \item" error.
  lbl/.style   = {font=\scriptsize\itshape, text=black!65, fill=white,
                  inner sep=2pt, align=center}
}

% forest default for directory trees
\forestset{
  dirtree/.style={
    for tree={
      font=\ttfamily\small, grow'=0, child anchor=west, parent anchor=south,
      anchor=west, calign=first, inner sep=2pt,
      edge path={\noexpand\path [draw, \forestoption{edge}]
        (!u.south west) +(9pt,0) |- (.child anchor) \forestoption{edge label};},
      before typesetting nodes={if n=1{insert before={[,phantom]}}{}},
      fit=band, before computing xy={l=15pt}
    }
  }
}

% ===== misc ==============================================================
\captionsetup{font=small, labelfont={bf,color=inkblue}}
\setlist[itemize]{leftmargin=1.2em, itemsep=2pt, topsep=3pt}
\setlist[enumerate]{leftmargin=1.4em, itemsep=2pt, topsep=3pt}
\newcommand{\file}[1]{\texttt{\small #1}}
\newcommand{\code}[1]{\texttt{\small #1}}
\setcounter{tocdepth}{2}
\setcounter{secnumdepth}{3}

%   \begin{document} ... \end{document}
%
% Build:  pdflatex -interaction=nonstopmode deep-dive.tex   (run it three times, so the
%         table of contents and the TikZ positioning settle)
%
% PACKAGE NOTES. Every package used here was checked present on the machine these
% documents were built on (Debian TeX Live 2023). If you rebuild elsewhere and a build
% fails, these are the likely gaps. Do not add a package without checking first:
%   kpsewhich <pkg>.sty
%
% Deliberately NOT used, because they were absent and the build dies without them:
%   algorithm2e, algorithm, algorithmic, algpseudocode  (texlive-science not installed)
%     -> the \begin{pseudocode} environment defined below stands in for them.
%   lmodern, charter, mathptmx, times, mathpazo, courier, newtx*, tgtermes
%     -> their .sty files RESOLVE but the font files themselves are absent, so the
%        build dies with "I can't find file `bchr8t'" or similar. Verified by
%        test-compiling each one. Only Computer Modern (the default) and helvet
%        actually render. kpsewhich reporting a font it cannot render is the trap:
%        do not "upgrade" the font stack without test-compiling it first.
%   amsmath, amssymb -> not loaded, so \text{} and \blacksquare are unavailable.
%
% A fuller font stack and algorithm2e would come from:
%   sudo apt install texlive-fonts-recommended texlive-science cm-super

\usepackage[T1]{fontenc}
\usepackage[utf8]{inputenc}
% Font: Computer Modern (default). See the note above before changing this.
\usepackage[a4paper,margin=2.4cm,headheight=14pt]{geometry}
% microtype WITHOUT font expansion. This is load-bearing, do not "restore" it:
% cm-super is not installed, so T1 Computer Modern resolves to BITMAP fonts, and
% pdfTeX's auto-expansion needs scalable ones. With expansion on, any document
% long enough to reach a second page dies with
%   "pdfTeX error (font expansion): auto expansion is only possible with scalable fonts"
% Protrusion (the bigger visual win) still works and stays on.
% `sudo apt install cm-super` would lift this; until then, leave expansion off.
\usepackage[protrusion=true,expansion=false]{microtype}
% Kill hyphen ligatures in TYPEWRITER text. This is load-bearing, not cosmetic. Under T1,
% \texttt{--verbose} silently renders as "-verbose", and \texttt{---} as "--", so a
% document can print a command-line flag as the WRONG flag and look perfectly fine doing
% it. Verified by rendering the page and reading the glyphs. lstlisting and \verb were
% always safe; this makes \code, \file and \texttt safe too, which is what prose uses.
% Ordinary prose is unaffected and still ligates, which is what you want there.
\DisableLigatures[-]{family=tt*}
\usepackage{graphicx}
\usepackage{xcolor}
\usepackage{booktabs}
\usepackage{enumitem}
\usepackage{fancyhdr}
\usepackage{titlesec}
\usepackage{tcolorbox}
\usepackage{listings}
\usepackage{float}
\usepackage{caption}
\usepackage{tikz}
\usepackage{forest}
\usepackage{pgfplots}
\pgfplotsset{compat=1.18}
\usetikzlibrary{arrows.meta,positioning,shapes.geometric,fit,backgrounds,calc,decorations.pathreplacing}
\tcbuselibrary{skins,breakable}
\usepackage[hidelinks,bookmarks=true,bookmarksnumbered=true]{hyperref}

% ===== palette ===========================================================
\definecolor{inkblue}{HTML}{1F3A5F}
\definecolor{accent}{HTML}{2E7DAF}
\definecolor{softgrey}{HTML}{F4F6F8}
\definecolor{linegrey}{HTML}{C8D0D8}
\definecolor{codebg}{HTML}{F7F7F4}
\definecolor{warnred}{HTML}{A33B3B}
\definecolor{okgreen}{HTML}{2E7D53}

% ===== headings ==========================================================
\titleformat{\section}{\Large\bfseries\color{inkblue}}{\thesection}{0.7em}{}
\titleformat{\subsection}{\large\bfseries\color{inkblue}}{\thesubsection}{0.6em}{}
\titleformat{\subsubsection}{\normalsize\bfseries\color{accent}}{\thesubsubsection}{0.5em}{}

\pagestyle{fancy}
\fancyhf{}
\fancyhead[L]{\small\color{inkblue}\nouppercase{\leftmark}}
\fancyhead[R]{\small\color{inkblue}\thepage}
\renewcommand{\headrulewidth}{0.4pt}
\renewcommand{\headrule}{\hbox to\headwidth{\color{linegrey}\leaders\hrule height \headrulewidth\hfill}}

% ===== code ==============================================================
\lstdefinestyle{house}{
  backgroundcolor=\color{codebg},
  basicstyle=\ttfamily\footnotesize,
  keywordstyle=\color{inkblue}\bfseries,
  commentstyle=\color{gray}\itshape,
  stringstyle=\color{okgreen},
  numbers=left, numberstyle=\tiny\color{gray}, numbersep=8pt,
  frame=single, rulecolor=\color{linegrey},
  breaklines=true, breakatwhitespace=false,
  showstringspaces=false, tabsize=2,
  captionpos=b, xleftmargin=14pt, framexleftmargin=14pt,
  columns=fullflexible, keepspaces=true,
  upquote=true,  % render ' and " as straight quotes, not curly
  % NON-ASCII INSIDE A LISTING IS A HARD BUILD FAILURE, not a cosmetic issue: listings
  % reads raw bytes, so inputenc dies on the first multi-byte character ("Invalid UTF-8
  % byte sequence") and the whole document fails to build. These mappings cover what
  % realistically turns up in source code: Latin-1 accents, Spanish punctuation, smart
  % quotes picked up from a word processor, the degree and section signs.
  % NOT COVERED, and it cannot be here: Urdu, Arabic and Balochi script. pdflatex has no
  % font for them. Describe or transliterate such text; never paste it into a listing.
  literate=%
    {á}{{\'a}}1 {é}{{\'e}}1 {í}{{\'i}}1 {ó}{{\'o}}1 {ú}{{\'u}}1
    {Á}{{\'A}}1 {É}{{\'E}}1 {Í}{{\'I}}1 {Ó}{{\'O}}1 {Ú}{{\'U}}1
    {à}{{\`a}}1 {è}{{\`e}}1 {ì}{{\`i}}1 {ò}{{\`o}}1 {ù}{{\`u}}1
    {ä}{{\"a}}1 {ë}{{\"e}}1 {ï}{{\"i}}1 {ö}{{\"o}}1 {ü}{{\"u}}1
    {Ä}{{\"A}}1 {Ö}{{\"O}}1 {Ü}{{\"U}}1 {ß}{{\ss}}1
    {ñ}{{\~n}}1 {Ñ}{{\~N}}1 {ç}{{\c c}}1 {Ç}{{\c C}}1
    {¿}{{\textquestiondown}}1 {¡}{{\textexclamdown}}1
    {°}{{\textdegree}}1 {§}{{\S}}1 {£}{{\pounds}}1 {€}{{EUR}}1
    {“}{{\textquotedblleft}}1 {”}{{\textquotedblright}}1
    {‘}{{\textquoteleft}}1 {’}{{\textquoteright}}1
    {–}{{-}}1 {—}{{-}}1 {…}{{...}}1 {→}{{$\rightarrow$}}1 {✓}{{x}}1
}
\lstset{style=house}

% ===== callout boxes =====================================================
% \begin{keyidea}{Title} ... \end{keyidea}   - a concept a beginner must grasp
\newtcolorbox{keyidea}[1]{
  enhanced, breakable, colback=softgrey, colframe=accent,
  boxrule=0.4pt, arc=2pt, left=8pt, right=8pt, top=6pt, bottom=6pt,
  title={\bfseries #1}, coltitle=white, colbacktitle=accent, fonttitle=\small
}

% \begin{jargon}{Term} definition \end{jargon}  - define a term at first use
\newtcolorbox{jargon}[1]{
  enhanced, breakable, colback=white, colframe=linegrey,
  boxrule=0.4pt, arc=2pt, left=8pt, right=8pt, top=6pt, bottom=6pt,
  title={\bfseries Jargon: #1}, coltitle=inkblue, colbacktitle=softgrey, fonttitle=\small
}

% \begin{honest}{Title} ... \end{honest}  - brutally honest finding (rule 18)
\newtcolorbox{honest}[1]{
  enhanced, breakable, colback=white, colframe=warnred,
  boxrule=0.6pt, arc=2pt, left=8pt, right=8pt, top=6pt, bottom=6pt,
  title={\bfseries #1}, coltitle=white, colbacktitle=warnred, fonttitle=\small
}

% \begin{pseudocode}{Caption} ... \end{pseudocode}
% Stands in for algorithm2e, which is NOT installed. Use \Step / \Indent inside.
\newtcolorbox{pseudocode}[1]{
  enhanced, breakable, colback=codebg, colframe=inkblue,
  boxrule=0.4pt, arc=2pt, left=10pt, right=8pt, top=6pt, bottom=6pt,
  fontupper=\ttfamily\small,
  title={\bfseries\rmfamily #1}, coltitle=white, colbacktitle=inkblue, fonttitle=\small
}
\newcommand{\Step}[1]{\par\noindent #1}
\newcommand{\Indent}[1]{\par\noindent\hspace*{2em}#1}

% ===== tikz house styles =================================================
\tikzset{
  box/.style   = {rectangle, rounded corners=2pt, draw=inkblue, fill=softgrey,
                  text=inkblue, align=center, inner sep=6pt, minimum height=9mm,
                  font=\small},
  extbox/.style= {rectangle, rounded corners=2pt, draw=linegrey, fill=white,
                  text=black!70, align=center, inner sep=6pt, minimum height=9mm,
                  font=\small, dashed},
  store/.style = {cylinder, shape border rotate=90, aspect=0.25, draw=accent,
                  fill=white, text=inkblue, align=center, inner sep=5pt,
                  font=\small},
  dec/.style   = {diamond, aspect=1.6, draw=accent, fill=white, text=inkblue,
                  align=center, inner sep=2pt, font=\scriptsize},
  flow/.style  = {-{Stealth[length=2.4mm]}, draw=inkblue, thick},
  dflow/.style = {-{Stealth[length=2.4mm]}, draw=accent, thick, dashed},
  % align=center is load-bearing: without it, a \\ inside an edge label throws
  % the misleading "perhaps a missing \item" error.
  lbl/.style   = {font=\scriptsize\itshape, text=black!65, fill=white,
                  inner sep=2pt, align=center}
}

% forest default for directory trees
\forestset{
  dirtree/.style={
    for tree={
      font=\ttfamily\small, grow'=0, child anchor=west, parent anchor=south,
      anchor=west, calign=first, inner sep=2pt,
      edge path={\noexpand\path [draw, \forestoption{edge}]
        (!u.south west) +(9pt,0) |- (.child anchor) \forestoption{edge label};},
      before typesetting nodes={if n=1{insert before={[,phantom]}}{}},
      fit=band, before computing xy={l=15pt}
    }
  }
}

% ===== misc ==============================================================
\captionsetup{font=small, labelfont={bf,color=inkblue}}
\setlist[itemize]{leftmargin=1.2em, itemsep=2pt, topsep=3pt}
\setlist[enumerate]{leftmargin=1.4em, itemsep=2pt, topsep=3pt}
\newcommand{\file}[1]{\texttt{\small #1}}
\newcommand{\code}[1]{\texttt{\small #1}}
\setcounter{tocdepth}{2}
\setcounter{secnumdepth}{3}


\title{\textbf{Trade Intelligence Copilot}\\[4pt]
\large The Brief: what it is, how it works, and what it actually scored}
\author{Explainer for Salman Adnan's portfolio}
\date{}

\begin{document}
\maketitle
\thispagestyle{empty}

% =====================================================================
\section{The one-minute version}

Ask a plain-English question about Pakistani sugar-sector import data and get a
number that came out of a real database, with the exact query shown. Ask about a
trade rule and get a definition with the source document cited. Ask about
something the data does not cover, such as Pakistan's wheat exports in 2019, and
get a refusal instead of an invention.

\begin{keyidea}{The headline is the evaluation, not the copilot}
Wiring a language model to a database is a weekend. Knowing whether it is
\emph{right} is the engineering. This project ships a 28-item test set whose
correct answers were computed directly from the database, runs the whole system
against the live Gemini API, publishes the real numbers including the failures,
and then checks its own automatic grader against independent labels (11 of 12
agree).
Everything else in this document exists to support that sentence.
\end{keyidea}

\begin{table}[H]
\centering
\small
\begin{tabular}{@{}ll@{}}
\toprule
\textbf{Metric} & \textbf{Result} \\
\midrule
SQL answer correctness, versus gold computed from the database & \textbf{15 / 15 = 100\%} \\
Retrieval hit rate & \textbf{8 / 10 = 80\%} \\
Refusal precision / recall on unanswerable questions & \textbf{1.00 / 1.00} \\
LLM-as-judge faithfulness (\code{gemini-2.5-pro}) & \textbf{21 / 22 = 95.5\%} \\
Judge versus independent labels (12 items) & \textbf{11 / 12 = 91.7\%} \\
Mean cost per query & \textbf{\$0.00093} \\
Latency p50 / p95 & \textbf{9.0 s / 27.0 s} (8.8 s / 12.5 s without 3 rate-limited items) \\
\bottomrule
\end{tabular}
\caption{From one real run over \file{eval/testset.jsonl}. Section~\ref{sec:verify}
states exactly which of these I re-derived myself and which rest on the README.}
\end{table}

% =====================================================================
\section{The four ideas you need}

\begin{jargon}{Retrieval-augmented generation (RAG)}
A language model knows only what it was trained on, and when it does not know it
invents fluent nonsense. RAG fixes this by fetching the relevant facts from a
trusted source first, then handing the model \emph{only} those facts and telling
it to use nothing else. Closed-book exam becomes open-book, with citations.
\end{jargon}

\begin{jargon}{Agentic retrieval}
Plain RAG always runs the same search. Here a model reads the question first and
\emph{decides} which sources it needs: the database, the document corpus, an
HS-code lookup, or several at once. That decision step is what makes it
``agentic''. The model picks the tools; ordinary code runs them.
\end{jargon}

\begin{jargon}{Text-to-SQL}
SQL is the language for querying tables of data. Text-to-SQL means the model
translates ``how many shipments came from China?'' into a query, and the
\emph{database} computes the number. The model is never asked to recall or
estimate a figure, which is the whole point: if the query is right, the number
is right by construction. The danger is obvious, since a model is now writing
executable code against your data, and Section~\ref{sec:safety} is about
containing that.
\end{jargon}

\begin{jargon}{Evaluation layer}
Machinery that measures whether the system is actually right on a fixed set of
questions with known answers, producing numbers you can publish and argue with.
Not a demo. A score. Crucially, the ``known answers'' here are not human
opinions: they are computed by querying the database directly, so ``correct''
is not a judgment call.
\end{jargon}

% =====================================================================
\section{How it works}

\begin{figure}[H]
\centering
\resizebox{\textwidth}{!}{
\begin{tikzpicture}[node distance=8mm]
  \node[box] (q) {Question};
  \node[box, right=12mm of q, fill=accent!12] (plan) {\textbf{Planner}\\\scriptsize gemini-2.5-flash\\\scriptsize picks tools, does\\\scriptsize not answer};

  \node[box, above right=7mm and 14mm of plan] (sql) {\textbf{text-to-SQL}\\\scriptsize LLM writes a SELECT};
  \node[box, right=14mm of plan] (ret) {\textbf{retrieval}\\\scriptsize embed + cosine\\\scriptsize over 48 chunks};
  \node[box, below right=7mm and 14mm of plan] (hs) {\textbf{hs\_lookup}\\\scriptsize plain SQL join};

  \node[box, right=12mm of sql, fill=warnred!12] (guard) {\textbf{SQL guard}\\\scriptsize + read-only DB};

  \node[dec, right=36mm of ret, xshift=10mm] (ev) {any\\evidence?};
  \node[box, above right=8mm and 12mm of ev, fill=warnred!12] (refuse) {\textbf{Refuse}\\\scriptsize costs nothing};
  \node[box, below right=8mm and 12mm of ev, fill=accent!12] (syn) {\textbf{Synthesizer}\\\scriptsize evidence only};
  \node[box, right=12mm of ev, xshift=28mm] (out) {answer + citations\\+ trace + cost};

  \draw[flow] (q) -- (plan);
  \draw[flow] (plan) -- (sql); \draw[flow] (plan) -- (ret); \draw[flow] (plan) -- (hs);
  \draw[flow] (sql) -- (guard);
  \draw[flow] (guard) -| (ev); \draw[flow] (ret) -- (ev); \draw[flow] (hs) -| (ev);
  \draw[flow] (ev) -- node[lbl,above]{no} (refuse);
  \draw[flow] (ev) -- node[lbl,below]{yes} (syn);
  \draw[flow] (refuse) -- (out); \draw[flow] (syn) -- (out);
\end{tikzpicture}}
\caption{Three model calls on a typical answer (plan, write SQL, synthesize);
two on a refusal, because the evidence check fires before the expensive call.}
\end{figure}

Read it as a sentence: \emph{an LLM decides which tools to run; ordinary code
runs them; a cheap check refuses when they return nothing; a second LLM writes
prose using only what they returned.}

\begin{keyidea}{No model is ever trusted with both the facts and the final word}
The planner chooses but does not answer. The text-to-SQL model writes a query
but never executes it. The synthesizer writes prose but is handed a fixed block
of evidence. Every model in the system is boxed in by code on both sides.
\end{keyidea}

\subsection{What is underneath}

\begin{itemize}
\item \textbf{The trade data}: 2,000 import shipments in a single SQLite file
  (a complete database that is just a file, no server), produced by a seeded
  generator, \file{data/build\_db.py}. Every company, price and quantity is made
  up; only the HS headings and country names are real. All imports, all HS
  chapter 17 (sugar), all between 2025-01-01 and 2025-12-31, from 28 countries.
  Those scope facts are written into the prompts, which is what lets the system
  know what it cannot answer.
\item \textbf{The regulations corpus}: 16 authored plain-language explainers,
  split at their section headings into 48 chunks, each embedded as 768 numbers
  representing its meaning, cached on disk.
\item \textbf{The vector store}: a 48-by-768 numpy matrix with pre-normalized
  rows, so finding the closest chunk is a single matrix-vector product. Exact,
  instant, 57 lines.
\item \textbf{No framework.} About a thousand lines of Python, no LangChain, no
  agent library. For a project whose purpose is to show the author understands
  the mechanism, a framework would hide the interesting parts.
\end{itemize}

\begin{keyidea}{The pgvector decision, defended}
pgvector was the first choice; the target machine has PostgreSQL 18 without the
\code{vector} extension and no way to install it. Rather than fight the
infrastructure, the author wrote the search by hand. For 48 chunks this is the
better engineering anyway: brute force is instant and \emph{exact}, whereas an
approximate index would be slower to build, approximate, and one more moving
part. The honest cost, no ANN index and no server-side filtering, starts to
matter around six figures of chunks, not 48.
\end{keyidea}

% =====================================================================
\section{Two things that make it trustworthy}

\subsection{The safety boundary}
\label{sec:safety}

A language model is writing SQL against a database, so the SQL is contained in
three independent layers:

\begin{enumerate}
\item \textbf{A guard} that rejects anything which is not a single
  \code{SELECT}: no comments, no second statement, no write or DDL keywords, no
  \code{PRAGMA}.
\item \textbf{The connection}, opened \code{mode=ro\&immutable=1} with
  \code{PRAGMA query\_only = ON}, so SQLite itself refuses writes.
\item \textbf{A row cap} of 200, so a runaway \code{SELECT} cannot flood memory
  or a prompt.
\end{enumerate}

The argument is not that any layer is strong. A keyword blocklist is a
blocklist and someone will think of a word you forgot. The argument is that a
bypass must defeat all three, and the tests assert each one independently,
including that a raw \code{CREATE TABLE} fails on the connection even with the
guard skipped entirely.

\subsection{The refusal path}

If every tool comes back empty, a nine-line function refuses immediately,
before any synthesis call. It cannot be argued out of it, cannot be
prompt-injected, costs zero tokens and zero seconds, and returns a fixed string.
Faster, cheaper and more reliable at the same time, which is a rare trade.

% =====================================================================
\section{The evaluation layer}

\begin{figure}[H]
\centering
\resizebox{0.96\textwidth}{!}{
\begin{tikzpicture}[node distance=8mm]
  \node[box] (ts) {\file{testset.jsonl}\\\scriptsize 28 items + gold};
  \node[box, right=12mm of ts, fill=accent!12] (ask) {\textbf{the real copilot}\\\scriptsize live API};
  \node[box, right=12mm of ask, fill=accent!12] (judge) {\textbf{judge}\\\scriptsize gemini-2.5-pro};
  \node[store, right=12mm of judge] (tr) {\file{transcripts.json}};

  \node[box, below=18mm of ask, fill=warnred!12] (metrics) {\textbf{compute\_metrics()}\\\scriptsize offline, free, repeatable};
  \node[store, left=14mm of metrics] (hl) {\file{human\_labels.json}\\\scriptsize 12 independent labels};
  \node[store, right=14mm of metrics] (res) {\file{results.json}\\\scriptsize published numbers};

  \draw[flow] (ts) -- (ask); \draw[flow] (ask) -- (judge); \draw[flow] (judge) -- (tr);
  \draw[dflow] (tr.south) -- ++(0,-7mm) -| node[lbl,pos=0.25,above]{saved once, scored forever} (metrics.north);
  \draw[flow] (hl) -- (metrics);
  \draw[flow] (metrics) -- (res);
\end{tikzpicture}}
\caption{Split in two on purpose. The paid run happens once and is saved; the
scoring can then be re-run, corrected and independently audited for free.}
\end{figure}

The 28 items are 12 SQL questions, 7 retrieval questions, 3 that need both, and
6 that must be refused. SQL correctness is scored over the 12 plus the 3, which
is where \textbf{15/15} comes from.

\begin{keyidea}{Ground truth is what makes the number mean anything}
The gold answer to ``how many shipments came from China?'' is not a human guess.
It is the result of running the counting query against the database directly.
That single decision is what caught the project's best bug, in its first version:
asked how many shipments one importer brought in, the model wrote a \code{buyer
LIKE} match on the full company name without wildcards, matched nothing, and
reported \emph{0 shipments}. The answer was fluent and completely wrong. The gold
number, which was dozens, exposed it in one run. Tightening the schema prompt to
require \code{LIKE '\%...\%'} for names fixed it, and the same kind of question on
the synthetic data (``How many shipments did Orvella Foods import?'') now comes
back right: 191.
\end{keyidea}

\begin{keyidea}{Faithful is not the same as correct}
That zero-row answer was judged \textbf{faithful}: it accurately reported what
its query returned. It was also completely wrong. This is why the project
measures grounding and ground-truth correctness as two separate numbers. Anyone
reporting a single ``accuracy'' figure for a RAG system has not hit this yet.
\end{keyidea}

\begin{keyidea}{A judge you have not validated is not evidence}
If you grade a model with a model, who grades the judge? The harness compares
the judge with independent labels on a 12-item subset, each answer read against
the evidence it was allowed to use. Labels belong to one run's answers, so after
the switch to synthetic data the subset was labelled afresh rather than carried
over: 11 of 12 agree, the same as on the first version. The single disagreement is
the judge's single miss: on ``how many shipments under HS code 1702.1100
(lactose)'', the answer had the right count (227) and repeated the question's own
``(lactose)'' label, which the evidence rows never state. The judge called that
unfaithful; the labels call it faithful, because its one load-bearing claim, the
count, is exactly the query result. Both are defensible, which is exactly the
point.
\end{keyidea}

% =====================================================================
\section{What I verified, and what I did not}
\label{sec:verify}

\begin{keyidea}{The live run was re-done for the synthetic data}
When the data was replaced, the live phase was run again once (28 items, about
\$0.026 of API usage), and the four live tests were run against the API. Every
other check below is offline.
\end{keyidea}

\begin{itemize}
\item \textbf{The gold answers are genuinely ground truth: verified.} They are
  written by \file{eval/build\_testset.py}, which runs a query for each one, and I
  re-derived them with independent queries (total USD 274,995,234.91; 2,000 rows;
  China 208 shipments; top country Brazil; Orvella Foods 191; HS 1704 total
  11,665,069.33). Every scope fact checked out too: \code{trade\_type} has exactly
  one distinct value, \code{'import'}; dates run 2025-01-01 to 2025-12-31; 28
  countries; 28 companies. Two builds of the database are byte-identical.
\item \textbf{The metrics reproduce: verified.} With credentials unset so no API
  call was possible, \code{python -m eval.run\_eval --metrics} recomputed every
  number from the committed transcripts and produced output
  \textbf{byte-identical} to the committed \file{results.json}.
\item \textbf{The test suite: verified.} \code{pytest} with no credentials gives
  \textbf{29 passed, 4 skipped}; with credentials the four live tests pass too,
  matching the README's claim exactly.
\end{itemize}

\begin{honest}{Precisely what ``15/15'' rests on}
All three links are verified for the current numbers: the gold answers are
provably correct (re-derived from the database), the scoring provably yields
15/15 (recomputed offline, byte-identical), and the transcripts are the output of
a real live run against Gemini, done on 2026-10-14 when the data was replaced.

That run hit Vertex rate limits. Three items waited out \code{429} responses, so
their recorded latencies (27, 46 and 30 seconds) include waiting, which is why
p95 is reported both with and without them. The transcripts carry per-call token
counts and genuine \emph{failures} (two retrieval misses and one judged
unfaithful answer); nothing was re-run to make a number look better.
\end{honest}

% =====================================================================
\section{Honest findings}

\subsection{What the project already owns, and gets right}

The wildcard bug and the faithful-but-wrong distinction above are both the
project's own findings, documented in its README. So is a third: the first eval
run flagged HS-lookup answers as unfaithful because the judge could not see the
HS rows the answerer had used. Fixing that plumbing turned the remaining flags
into genuine findings. A harness that manufactures failures is as dangerous as
one that hides them.

The two remaining faithfulness misses are a single pattern, and it is the best
technical observation in the project:

\begin{keyidea}{Enumeration is where grounded models leak}
In the first version, point-fact answers stayed grounded across the whole set,
and ``list everything under this heading'' is where the model reached past its
evidence, padding its list with plausible product names the judge caught by name.
A list of three items \emph{feels} incomplete, and the model's training pushes it
toward a satisfying list. In the synthetic re-run the lists stayed inside the
evidence, most likely because the synthetic HS hierarchy is small enough for the
lookup to return every item. That is luck, not a fix. The fix is structural, not a
sterner instruction: post-filter enumerated items against the evidence rows.
\end{keyidea}

\subsection{What this review found}

\begin{honest}{Five things worth knowing before an interviewer finds them}
\begin{enumerate}
\item \textbf{\code{in\_scope} is dead code.} The planner is prompted for it, the
  model reasons about it, the code carefully normalizes it, and \emph{nothing
  ever reads it}. It costs tokens on every request and reads like a safety
  control while being none. Either wire it in (it could skip straight to refusal
  and save two model calls) or delete it.
\item \textbf{The SQL guard rejects valid queries.} The keyword regex scans the
  whole query with no awareness of string literals, so
  \code{... WHERE name LIKE '\%Create\%'} is rejected. I ran this against the
  real guard to confirm it. A company named ``Create Foods'' would be
  unqueryable. It fails \emph{safe}, which is the right direction to err, but
  the clean fix is SQLite's authorizer callback, which cannot be fooled by
  quoting.
\item \textbf{``Deterministic refusal'' is deterministic only in its last step.}
  All six refusals traced as \code{['plan', 'refusal']}, meaning no tool ran at
  all. The guard did not weigh evidence and find it wanting; it fired because
  the \emph{planner} chose no tools. The decisive judgment belongs to an LLM.
  6/6 is real, but the honest claim is narrower than ``enforced
  deterministically''.
\item \textbf{The 80\% retrieval hit rate is really a planner score.} Both
  misses had no retrieval step in the trace at all: the planner never invoked
  retrieval, so no chunk could be found. Retrieval went \textbf{8 for 8} on
  everything it was actually asked to find. The metric's name points at the
  wrong component.
\item \textbf{A latent citation bug.} The SQL rows shown to the judge are gated
  on \code{used\_sql}, a field the \emph{synthesizer reports about itself},
  rather than on whether rows exist. That is a structural repeat of the
  HS-lookup bug already fixed once. I checked all 28 transcripts and found
  \textbf{zero} instances, so the published 95.5\% is not contaminated. It is a
  latent fragility with a one-line fix.
\end{enumerate}
\end{honest}

Minor: \code{jinja2} is pinned in \file{requirements.txt} and never imported
(the server reads its template as a plain string). Per-query costs are real
token counts multiplied by \emph{list prices hardcoded in a dict}, not billed
amounts, so they will silently drift if Google changes pricing.

\subsection{Limits the project states, and should keep stating}

\begin{itemize}
\item \textbf{One run, 28 items.} Every number is a single sample. 15/15 on 15
  questions is a good result with wide error bars. ``The paths work'', not ``it
  is 100\% accurate''.
\item \textbf{SQL correctness is scraped from prose}, so a wrong answer
  containing the right number would score correct. The project documents this
  itself.
\item \textbf{The judge and the answerer are the same model family}, so shared
  blind spots are possible; the independent labels only partly cover this.
\item \textbf{The corpus is authored summaries, not primary law.} Faithfulness
  measures grounding in that corpus, not legal accuracy. The project says so in
  three places and should keep doing so.
\item \textbf{Latency is high} (p50 about 9 s) because three or four model calls
  run strictly in sequence, even when the tools are independent.
\end{itemize}

% =====================================================================
\section{How to defend it}

\begin{keyidea}{The five things to say}
\begin{enumerate}
\item \textbf{Lead with the eval.} Anyone can wire an LLM to a database. The
  differentiator is gold answers computed from the database, published metrics
  including the failures, and a judge checked against independent labels.
\item \textbf{Tell the wildcard story.} A fluent answer said 0 shipments; the
  truth was dozens; the gold number caught it in one run. That is the argument
  for evaluation in thirty seconds.
\item \textbf{Separate faithful from correct}, and note the same bug proved why:
  the wrong answer was perfectly faithful to its own wrong query.
\item \textbf{Defend the hand-built vector store.} 48 chunks, exact search, one
  matrix-vector product, no server. The simpler thing is also the better thing
  here.
\item \textbf{Volunteer the limits.} 80\% is really a routing problem.
  Enumeration is where grounding leaks. \code{in\_scope} is dead. Say it first.
\end{enumerate}
\end{keyidea}

\subsection{Expect these questions}

\begin{itemize}
\item \emph{``What stops the model dropping your tables?''} Three layers, a test
  asserts a write fails even with the guard bypassed. Then volunteer the
  string-literal false positive.
\item \emph{``Isn't using an LLM to grade an LLM circular?''} Yes, which is why
  the judge is checked against independent labels: 11/12, on both the first
  version and the synthetic re-run, labelled afresh each time.
\item \emph{``Faithfulness is 95.5\%; why should I trust it?''} Because the judge
  agrees with independent labels 11 times in 12, and its one miss is a judgment
  call about a label the user supplied. On the first version it was 86\%, and the
  misses were enumeration padding, a known failure mode with a structural fix.
\item \emph{``Would this scale?''} The vector store needs a real index somewhere
  around six figures of chunks. The tighter constraint is latency: merge the
  planner with text-to-SQL and parallelise independent tools.
\end{itemize}

% =====================================================================
\section{Glossary}

\begin{description}[leftmargin=!,labelwidth=3.2cm,style=nextline]
\item[Agentic retrieval] Retrieval where a model decides which sources to
  consult rather than always running one fixed search.
\item[Chunk] One retrievable piece of a document; here, one section.
\item[Cosine similarity] How closely two vectors point the same way. With
  normalized vectors it is just the dot product.
\item[Embedding] A list of numbers (768 here) representing the meaning of text.
\item[Faithfulness] Whether an answer is supported by the evidence it was given.
  Not the same as correctness.
\item[Ground truth] The known-correct answer, established independently; here,
  computed straight from the database.
\item[HS code] International product classification. This dataset is chapter 17,
  sugar.
\item[LLM-as-judge] Using a model to grade another model's output.
\item[Precision / recall] Of the refusals, how many were right / of the questions
  needing refusal, how many were caught.
\item[RAG] Fetch the facts, then let the model write using only those facts.
\item[SQLite] A complete relational database that is a single file, no server.
\item[Text-to-SQL] Having a model translate a question into a real database
  query and running it.
\end{description}

\vfill
\begin{center}
\small\itshape
Built from the source at \file{trade-intelligence-copilot}. Every metric quoted
was recomputed offline from the committed transcripts, and every gold answer
re-derived from the shipped database, before it was written down.
\end{center}

\end{document}
